\section{Research Trends and Emerging Directions}
\label{sec:trends}

Sections~\ref{sec:failure} and~\ref{sec:methodology} document what tends
to go wrong when modern frameworks enter the GrC pipeline.  This section
takes a constructive view: three research directions are showing sustained
momentum and---when pursued carefully---genuine potential.  We assess each
against the Pattern~A/B/C diagnostic and identify the precise conditions
under which a contribution would be non-trivial.

%-----------------------------------------------------------------------
\subsection{Incremental and Stream-Aware Feature Selection}
\label{sec:trend_incremental}
%-----------------------------------------------------------------------

The dominant evaluation setting in the GrC literature---a static decision
table of a few hundred objects---does not reflect the conditions under which
rough-set feature selection is most useful in practice.  Real data pipelines
in network intrusion detection, sensor fusion, and clinical monitoring produce
objects continuously; re-running a full forward-greedy reduct search after
each batch is infeasible.

Two recent lines address this directly.
Zhang et al.~\cite{Zhang2023TKDEincremental} (TKDE~2023, 50 citations)
develop an incremental GB-NRS algorithm that tracks which granular balls
are affected by an arriving batch $\Delta\mathcal{U}$ and updates attribute
significance locally, achieving $O(|\Delta\mathcal{U}| \cdot |\mathcal{U}|)$
update cost versus $O(|\mathcal{U}|^2 d)$ for a full re-run---a factor of
roughly $|\mathcal{U}|/d$ speedup.  Qian et al.~\cite{Qian2023InfoFusion}
(Information Fusion~2023, 46 citations) extend a similar idea to
partial-label settings, where arriving objects may carry uncertain class
memberships.

The algorithmic contribution---maintaining granular-ball membership lists
across updates---is real and practically useful.  The Pattern~C risk, noted
in Section~\ref{sec:geo_emerging}, is that evaluations report only time, not
reduct quality.  A non-trivial theoretical contribution would provide a bound
of the form
\[
  \Pr\bigl[R_{t+1} \neq R^*_{t+1}\bigr] \leq f\!\left(\frac{|\Delta\mathcal{U}|}{|\mathcal{U}|},\,
  \gamma(R^*_t),\, \tau\right),
\]
where $R^*_{t+1}$ is the globally optimal reduct on
$\mathcal{U}\cup\Delta\mathcal{U}$ and $\gamma(R^*_t)$ is the dependency
degree of the current reduct.  Intuitively: if the current reduct has a large
dependency margin and the arriving batch is small, a local update is likely to
remain globally optimal.  Formalising this in terms of the neighbourhood
graph's stability under batch insertion is an open problem
(Section~\ref{sec:open}).

\paragraph{Conditions for a non-trivial contribution.}
A new incremental GrC paper should (a) bound reduct deviation, not only report
speedup; (b) include a full-batch re-run baseline on accuracy (not just time);
(c) test on a genuinely streaming benchmark with at least 10 update rounds.

%-----------------------------------------------------------------------
\subsection{Three-Way Decision as a Payoff-Function Layer}
\label{sec:trend_3wd}
%-----------------------------------------------------------------------

Three-way decision (3WD) adds a \emph{deferred} option to the binary
positive/negative split, partitioning $\mathcal{U}$ into $\mathrm{POS}$
(accept), $\mathrm{NEG}$ (reject), and $\mathrm{BND}$ (boundary, withhold
decision).  Introduced by Yao~\cite{Yao2010threeway} as a cost-minimising Bayesian
partition, 3WD is now one of the most-cited paradigms in the rough set
community (1321 citations, retrieved September 2025).

3WD does not change the shape of the granule---it changes the payoff function
applied after the granule determines region membership.  Under a cost
matrix $(\lambda_{\mathrm{PP}}, \lambda_{\mathrm{NP}}, \lambda_{\mathrm{BP}},
\lambda_{\mathrm{PN}}, \lambda_{\mathrm{NN}}, \lambda_{\mathrm{BN}})$,
the expected risk of assigning an object to each region is computed, and the
region with minimum risk is chosen.  Yang et al.~\cite{Xia2024TFS3WC} combine
this with GBNRS: the granule's purity score provides the posterior estimate
required by the 3WD risk calculation.  The integration is natural and yields
a principled cost-sensitive feature selection pipeline.

The value of 3WD-GrC is concentrated in domains where misclassification costs are
asymmetric and known: medical screening (missing a positive case costs far
more than a false alarm), network intrusion detection, rare-event prediction.
In these settings, reporting the boundary region size alongside the cost-sensitive
accuracy provides a decision-support function that standard NRS cannot replicate.

\paragraph{Conditions for a non-trivial contribution.}
Two patterns must be checked before claiming a 3WD-GrC method outperforms
standard NRS.
\emph{Pattern~A}: at symmetric costs and $\tau \to 1$, the 3WD regions reduce
to the NRS positive/negative/boundary regions---the method adds no information
beyond what standard NRS already provides.  The cost matrix must be
demonstrably non-symmetric, and sensitivity to cost-matrix choice must be
reported.
\emph{Pattern~C}: evaluating solely on cost-sensitive accuracy (which 3WD
directly optimises) while ignoring 0-1 accuracy and boundary region size
creates a self-confirming benchmark.  A fair evaluation reports all three
metrics: cost-sensitive accuracy, standard 0-1 accuracy on classified objects,
and the fraction of objects assigned to $\mathrm{BND}$ (a large boundary
region is a cost, not a benefit).

One question remains open beneath all of this.
The cost matrix is currently chosen by the experimenter, not derived from
data.  Connecting the optimal cost matrix to the dataset's class-overlap
profile~\cite{Santos2022overlap} and the Bayes error estimate would make
3WD-GrC self-configuring and would close the Pattern~A gap at the same time.
We list this as an open problem in Section~\ref{sec:open}.

%-----------------------------------------------------------------------
\subsection{Neural-Augmented Granular Computing}
\label{sec:trend_neural}
%-----------------------------------------------------------------------

A third trend feeds learned representations into the GrC pipeline: instead
of operating on raw feature vectors, the granule is defined in a latent
space produced by a neural encoder.
Zhang et al.~\cite{Zhang2025NN} use an ANN to assign per-feature scale
weights before computing granular ball purity, allowing the scale of each
attribute to vary across the feature space.
Xie et al.~\cite{Xie2024TPAMI} (MGNR, TPAMI~2024, 57 citations) extract
multi-granularity neighbourhood representations across multiple scales
simultaneously, achieving state-of-the-art accuracy on large benchmarks.

If the neural encoder is trained to separate classes, the latent granule
becomes semantically meaningful rather than geometrically arbitrary.  This
directly addresses Pattern~B: a well-trained encoder maps features to a space
where Euclidean distance reflects class structure, breaking the correlation
between a feature's scale and its granule-based importance score.

The risk persists, however, if the encoder is trained unsupervised (e.g., an
autoencoder) or if the latent representation is not verified to be
class-discriminative: in that case, the granule in latent space still tracks
magnitude properties of the learned representation rather than class structure.
Any neural-GrC paper must verify that the latent distance is not dominated by
the largest-variance latent dimension---the same diagnostic applied in
Section~\ref{sec:patB} for PAI, applied now to the neural representation.

\paragraph{Conditions for a non-trivial contribution.}
A neural-GrC paper should (a) compare against GBNRS on the same
benchmark with the same number of selected features; (b) verify that the
latent space is class-discriminative (e.g., silhouette score or linear
separability); (c) ablate the encoder: does removing the neural component
(reverting to raw features) substantially degrade results, or does the gain
come from the granule geometry rather than the representation?
